% predicate-logic-syntax.tex

%%%%%%%%%%%%%%%%%%%%
\begin{frame}{}
  \begin{center}
    {\Large ``先讲语法''}

    \vspace{0.30cm}
    \fig{width = 0.60\textwidth}{figs/syntax-semantics}
  \end{center}
\end{frame}
%%%%%%%%%%%%%%%%%%%%

%%%%%%%%%%%%%%%%%%%%
\begin{frame}{}
  \begin{definition}[一阶谓词逻辑的语言 (Language)]
    一阶谓词逻辑的语言包括以下7部分: \\[5pt]
    \begin{description}
      \setlength{\itemsep}{5pt}
      \item [逻辑联词:] $\lnot, \land, \lor, \to, \leftrightarrow$
      \item [\purple{量词符号:}] $\forall$ (\purple{forall}; 全称量词),
        $\exists$ (\purple{exists}; 存在量词)
      \item [变元符号:] $x, y, z, \dots$
      \item [左右括号:] $(, )$
      \pause
      \vspace{0.50cm}
      \item [\cyan{常数符号:}] 零个或多个常数符号 $a, b, c, \dots$, \teal{表达特殊的具体个体}
      \item [\blue{函数符号:}] $n$-元函数符号 $f, g, h, \dots$ ($n \in \blue{\N^{+}}$), \teal{表达个体上的运算}
      \item [\red{谓词符号:}] $n$-元谓词符号 $P, Q, R, \dots$ ($n \in \N$), \teal{表达个体的性质与关系}
    \end{description}
  \end{definition}

  \pause
  \[
    \text{初等数论的语言:} \qquad
    L = \set{\textbf{\cyan{0}}, \blue{S}, \blue{+}, \blue{\times}, \red{<}, \red{=}}
  \]
  \pause
  \vspace{-0.30cm}
  \[
    \text{集合论的语言:} \qquad
    L = \set{\red{=}, \red{\in}}
  \]
\end{frame}
%%%%%%%%%%%%%%%%%%%%

%%%%%%%%%%%%%%%%%%%%
\begin{frame}{}
  \begin{center}
    ``项''是一阶谓词逻辑要讨论的\red{个体}对象, \\
    用以构成\purple{公式} \gray{(它本身\blue{无所谓真假})}
  \end{center}

  \pause
  \begin{definition}[项 (Term)]
    \begin{enumerate}[(1)]
      \setlength{\itemsep}{8pt}
      \item 每个变元 $x, y, z, \dots$ 都是一个项;
      \item 每个常数符号都是一个项;
      \item 如果 $t_{1}, t_{2}, \dots, t_{n}$ 是项,
        且 \blue{$f$} 为一个 $n$-元\blue{\bf 函数}符号, \\
        则 $f(t_{1}, t_{2}, \dots, t_{n})$ 也是项;
      \item \teal{除此之外, 别无其它。}
    \end{enumerate}
  \end{definition}
  \[
    L = \set{\textbf{\cyan{0}}, \blue{S}, \blue{+}, \blue{\times}, \red{<}, \red{=}}
  \]
  \[
    x \qquad \cyan{\bf 0}
  \]
  \[
    S{\bf 0} \qquad x + SSS{\bf 0} \qquad (x + SSS{\bf 0}) \times y
    \pause \qquad {\bf 0} < S{\bf 0}
  \]
\end{frame}
%%%%%%%%%%%%%%%%%%%%

%%%%%%%%%%%%%%%%%%%%
\begin{frame}{}
  \begin{center}
    公式刻画了\red{个体的性质}或者\red{个体之间的关系}
    \gray{(它的语义就是它的\blue{真假值})}
  \end{center}

  \begin{definition}[\gray{合式 (Well-formed)} 公式 (Formula)]
    \begin{enumerate}[(1)]
      \setlength{\itemsep}{8pt}
      \item 如果 $t_{1}, \dots, t_{n}$ 是项, 且 $P$ 是一个 $n$ 元\blue{\bf 谓词}符号, \\
        则 $P(t_{1}, \dots, t_{n})$ 为公式, 称为\teal{\bf 原子公式};
      \item 如果 $\alpha$ 与 $\beta$ 都是公式,
        则 $(\alpha)$、$(\lnot \alpha)$、$(\alpha \;\red{\ast}\; \beta)$ 都是公式;
      \item 如果 $\alpha$ 是公式, $x$ 是一个变元, 则 \red{$\forall x\; \alpha$}
        与 \red{$\exists x\; \alpha$} 也是公式;
      \item \teal{除此之外, 别无其它。}
    \end{enumerate}
  \end{definition}

  \pause
  \vspace{0.20cm}
  \begin{center}
    \blue{\bf 约定:} 量词符号 \red{$\forall$} 与 \red{$\exists$} 的管辖范围尽可能短
    \pause
    \[
      \forall x\; \alpha \to \beta:
        \text{表示}\; (\forall x\; \alpha) \to \beta
        \quad \gray{\text{不表示}\; \forall x\; (\alpha \to \beta)}
    \]
    \pause
    \vspace{-0.30cm}
    \[
      \blue{\forall x.\; \alpha \to \beta}
    \]
  \end{center}
\end{frame}
%%%%%%%%%%%%%%%%%%%%

%%%%%%%%%%%%%%%%%%%%
\begin{frame}{}
  \[
    \ndnoname{\text{每个人都是要死的 \qquad 苏格拉底是人}}{\text{苏格拉底是要死的}}
  \]

  \pause
  \vspace{0.60cm}
  \[
    \ndnoname{\forall x\; (H(x) \to M(x)) \qquad H(s)}{M(s)}
  \]
\end{frame}
%%%%%%%%%%%%%%%%%%%%

%%%%%%%%%%%%%%%%%%%%
\begin{frame}{}
  \[
    L = \set{\textbf{\cyan{0}}, \blue{S}, \blue{+}, \blue{\times}, \red{<}, \red{=}}
  \]

  \vspace{0.30cm}
  \begin{enumerate}[(1)]
    \setlength{\itemsep}{8pt}
    \pause
    \item 0 不是任何自然数的后继
      \pause
      \[
        \forall x.\; \lnot (Sx = 0)
      \]
    \pause
    \vspace{-0.50cm}
    \item 两个自然数相等当且仅当它们的后继相等
      \pause
      \[
        \forall x.\; \forall y.\; (x = y \leftrightarrow Sx = Sy)
      \]
    \pause
    \vspace{-0.50cm}
    \item $x$ 是素数 ($x > 1$ 且 $x$ 没有除自身和1之外的因子)
      \pause
      \[
        \text{Prime}(x): S0 < x \land \forall y.\; \forall z.\; (y < x \land z < x) \to \lnot (y \times z = x)
      \]
    \pause
    \vspace{-0.50cm}
    \item 哥德巴赫猜想 (任一大于2的偶数, 都可表示成两个素数之和)
      \pause
      \begin{align*}
        &\forall x.\; (SS0 < 2 \land (\exists y.\; 2 \times y = x)) \to \\
          &\quad (\exists x_{1}.\; \exists x_{2}.\; \text{Prime}(x_{1}) \land \text{Prime}(x_{2})
            \land x_{1} + x_{2} = x)
      \end{align*}
  \end{enumerate}
\end{frame}
%%%%%%%%%%%%%%%%%%%%

%%%%%%%%%%%%%%%%%%%%
\begin{frame}{}
  \[
    \blue{\boxed{\lim_{x \to a} f(x) = b}}
  \]

  \pause
  \vspace{0.30cm}
  \begin{center}
    对于任意的正实数 $\epsilon$, 存在一个正实数 $\delta$, \\[6pt]
    使得对于任意的 $x$, 当 $0 < |x - a| < \delta$ 时, 都有 $|f(x) - b| < \epsilon$ 成立。
  \end{center}

  \pause
  \[
    \forall \epsilon \in \R^{+}.\; \exists \delta \in \R^{+}.\;
      \forall x \in \R^{+}.\; (0 < |x - a| < \delta \to |f(x) - b| < \epsilon)
  \]
\end{frame}
%%%%%%%%%%%%%%%%%%%%

%%%%%%%%%%%%%%%%%%%%
\begin{frame}{}
  \[
    \blue{\forall x \in A.\; \alpha} \quad\text{实际上是}\quad
      \blue{\forall x.\; (x \in A \;\red{\to}\; \alpha)} \quad\text{的简记}
  \]

  \uncover<2->{
    \fig{width = 0.45\textwidth}{figs/why-why-why}
  }

  \[
    \blue{\exists x \in A.\; \alpha} \quad\text{实际上是}\quad
      \blue{\exists x.\; (x \in A \;\red{\land}\; \alpha)} \quad\text{的简记}
  \]
\end{frame}
%%%%%%%%%%%%%%%%%%%%

%%%%%%%%%%%%%%%%%%%%
\begin{frame}{}
  motivation for substitution; and then for bound/free variables
  \[
    \forall x\; \phi \to \phi[t/x]
  \]
  \fig{width = 0.30\textwidth}{figs/substitution-red}

  \pause
  \[
    \phi = \exists y\; x \neq y
  \]

  \pause
  \[
    \blue{t = z:} \forall x\; \exists y\; x \neq y \to \exists y \blue{z} \neq y
  \]

  \pause
  \[
    \red{t = y:} \forall x\; \exists y\; x \neq y \to \exists y \red{y} \neq y
  \]
\end{frame}
%%%%%%%%%%%%%%%%%%%%

%%%%%%%%%%%%%%%%%%%%
\begin{frame}{}
  \begin{center}
    以下概念与程序设计语言中相应概念类似
  \end{center}
  \begin{columns}
    \column{0.50\textwidth}
      \fig{width = 0.80\textwidth}{figs/bound-free}
    \column{0.50\textwidth}
      \fig{width = 0.80\textwidth}{figs/integration}
  \end{columns}
\end{frame}
%%%%%%%%%%%%%%%%%%%%

%%%%%%%%%%%%%%%%%%%%
\begin{frame}{}
  \begin{exampleblock}{作用域 (Scope)、约束变元 (Bound)、自由变元 (Free)}
    \begin{enumerate}[(1)]
      \setlength{\itemsep}{6pt}
      \item $\forall \red{x}\; (P(\red{x}) \to Q(\red{x}))$
      \item $(\forall \red{x}\; P(\red{x})) \to Q(\green{x})$
      \item $\forall \red{x}\; \Big(P(\red{x}) \to \big(\exists \brown{y}\; R(\red{x}, \brown{y})\big)\Big)$
      \item $\Big(\forall \red{x}\; \forall \brown{y}\; \big(P(\red{x}, \brown{y})
        \land Q(\brown{y}, \green{z})\big)\Big) \land
        \exists \violet{x}\; P(\violet{x}, \green{y})$
    \end{enumerate}
  \end{exampleblock}
\end{frame}
%%%%%%%%%%%%%%%%%%%%

%%%%%%%%%%%%%%%%%%%%
\begin{frame}{}
  \begin{definition}[变元 $x$ 在公式 $\phi$ 中\green{\bf 自由出现}]
    $x$ 在 $\phi$ 中\green{\bf 自由出现}当且仅当
    \begin{description}[$\phi$ 是原子公式:]
      \item[$\phi$ 是原子公式:]
        $x$ 在 $\phi$ \purple{\bf 出现}。
      \item[$\phi = \lnot\psi$:]
        $x$ 在 $\phi$ {自由出现}。
      \item[$\phi = \psi \lor \gamma$:]
        $x$ 在 $\phi$ {自由出现}
        或 $x$ 在 $\gamma$ 中自由出现。
      \item[$\phi = \forall v\; \psi$:]
        $x$ 在 $\phi$ \teal{\bf 自由出现} 且 \teal{$x \neq v$}。
    \end{description}
  \end{definition}

  \[
    \Big(\forall \red{x}\; \forall \brown{y}\; \big(P(\red{x}, \brown{y})
      \land Q(\brown{y}, \green{z})\big)\Big) \land
      \exists \violet{x}\; P(\violet{x}, \green{y})
  \]

  \pause
  \vspace{0.50cm}
  \begin{definition}[变元 $x$ 在公式 $\phi$ 中\red{\bf 约束出现}]
    在 $\phi$ 中出现的变元如果不是自由出现, 则被称为\red{\bf 约束出现}。
  \end{definition}
\end{frame}
%%%%%%%%%%%%%%%%%%%%

%%%%%%%%%%%%%%%%%%%%
\begin{frame}{}
  \begin{definition}[改名 (Rename)]
    为尽量避免重名, 可将\red{\bf 约束变元}\teal{\bf 改名}为\blue{新鲜(fresh)}变元
  \end{definition}

  \[
    \Big(\forall x.\; \forall y.\; \big(P(x, y) \land Q(y, z)\big)\Big)
      \land \exists \blue{t}.\; P(\blue{t}, \teal{v})
  \]
\end{frame}
%%%%%%%%%%%%%%%%%%%%

%%%%%%%%%%%%%%%%%%%%
\begin{frame}{}
\end{frame}
%%%%%%%%%%%%%%%%%%%%

%%%%%%%%%%%%%%%%%%%%
\begin{frame}{}
  \begin{definition}[$t$ is free for $x$ in $\alpha$]
    \[
      y - 1 \text{ is \blue{free for} } x \text{ in } \exists z.\; (z < x)
    \]

    \[
      y - 1 \text{ is \red{{\it not} free for} } x \text{ in } \exists y.\; (y < x)
    \]
  \end{definition}

  \vspace{0.50cm}
  \begin{center}
    在公式 $\alpha$ 中, 项 $t$ 可以替换变量 $x$ (\blue{$\alpha[t/x]$})
  \end{center}
\end{frame}
%%%%%%%%%%%%%%%%%%%%